\section{Physical Explanations for the Accelerated Universe}

Since the discovery of the accelerated expansion of the Universe by 
\cite{Riess_1998} and \cite{Perlmutter_1999}, understanding its underlying origin has 
become one of the major challenges in modern cosmology. One possible explanation is the 
introduction of a cosmological constant $\Lambda$ in Einstein's equations. This constant 
can be interpreted in different ways: as a geometrical contribution modifying the Einstein 
equations, as an effective fluid with a constant energy density and equation of state 
$w=-1$, or as a manifestation of the vacuum energy of quantum fields. When interpreted 
as an additional energy component driving the accelerated expansion, it is referred to as 
dark energy. In the following subsections, we discuss several theoretical models proposed 
to explain the observed accelerated expansion.

\subsection{Cosmological Constant} \label{sec:expanding_universe}

The simplest way to account for the acceleration as we have mentioned before is by adding a cosmological constant $\Lambda$ to the Einstein equations \ref{eqn:einstein_lambda}, 
which is what is done in the $\Lambda$CDM model. We can look at $\Lambda$ in two equivalent ways, either a modification to the geometry of spacetime or as a part of the 
energy content of the Universe. Mathematically they are equivalent since the energy-momentum tensor $T_{\mu \nu}$ can include the $\Lambda g_{\mu \nu}$ and their 
equations of state are indistinguishable in $\Lambda$CDM. \\

\subsection{Dynamical Dark Energy}

One could naturally assume that if the acceleration of the Universe is driven by a fluid, its energy density would change throughout time. The common way to model this 
is with the Chevallier-Polarski-Linder (CPL) model (\cite{Chevallier_2001}) characterized by two parameters $w_0$ and $w_a$ such that:
\begin{equation}
    w(z) = w_0 + w_a \frac{z}{1+z} = w_0 + w_a(1-a)
\end{equation}
Therefore the Hubble parameter becomes:
\begin{equation}
    H(z) = H_0 \sqrt{\Omega_m(1+z)^3 + \Omega_r(1+z)^4 + \Omega_k(1+z)^2 + e^{\left( -3w_a\frac{z}{1+z} \right)} \Omega_{\text{DE}}(1+z)^{3(1+w_0+w_a)} }
\end{equation}
where $\Omega_{\text{DE}}$ is the energy density of this dynamical dark energy. $\Lambda$CDM then becomes a particular case where $w_0=-1$ and $w_a=0$.

\subsection{Quintessence Field}

The quintessence model uses an approach similar to inflation where one assumes the accelerated expansion is driven by a scalar field $\phi$, here called the quintessence field. 
The action associated to that field is written as:
\begin{equation}
    S(\phi) = \int d^4x \sqrt{-g}\left( -\frac{1}{2} g^{\mu \nu} \partial_{\mu}\phi \partial_{\nu} \phi - V(\phi) \right)
\end{equation}
The equations of motion of $\phi$ are determined by the Klein-Gordon equation:
\begin{equation}
    \ddot{\phi} + 3H \dot{\phi} + \frac{d}{d\phi}V = 0
\end{equation}
For a FLRW metric, the pressure and energy density become:
\begin{equation}
    P_{\phi} = \frac{1}{2}\dot{\phi}^2 - V(\phi)
\end{equation}
\begin{equation}
    \rho_{\phi} = \frac{1}{2}\dot{\phi}^2 + V(\phi)
\end{equation}
This gives the following equation of state:
\begin{equation}
    w = \frac{P_{\phi}}{\rho_{\phi}} = \frac{\frac{1}{2}\dot{\phi}^2 - V(\phi)}{\frac{1}{2}\dot{\phi}^2 + V(\phi)}
\end{equation}
Under a slow-roll approximation, the field evolves slowly such that $\dot{\phi} \simeq 0$ giving us $w\approx -1$ where we recover the results 
of a cosmological constant. If we consider a fast rolling field where the potential is negligible in front of the evolution, the equation of state becomes $w\approx 1$ 
however this is not possible since we know that for an accelerating Universe $w<-\frac{1}{3}$. Results from \cite{DESI_DR2_2025} show a preference for a 
dynamical dark energy in the CPL parameterization. \\ 

\subsection{Modified Theories of Gravity} \label{sec:modified_gravity}
An alternative to dark energy is to modify the theory of GR on cosmological scales. We call the models that do so Modified Gravity models (MG). Multiple classes of MG theories 
exist each modifying a certain aspect of GR. We will be presenting some of those models in each class.\\

\noindent \textbf{\underline{Additional scalar fields}} \\

The first class consists of adding an additional scalar field to the Einstein-Hilbert action. These fields need to have equations of motions of second-order time derivatives 
at most to avoid any unphysical solutions. \\

Models in this class include the $f(R)$ models (\cite{Sotiriou_2010}) which replaces the Ricci scalar $R$ by a general function of $R$. The action for an $f(R)$ gravity becomes:
\begin{equation}
    S = \frac{c^4}{16\pi G}\int \sqrt{-g} f(R) d^4x + \int \sqrt{-g} \mathcal{L}_m d^4x
\end{equation}
And the field equations are:
\begin{equation}
    f'(R)R_{\mu \nu} - \frac{1}{2}f(R)g_{\mu \nu} + [g_{\mu \nu}\Box -\nabla_{\mu}\nabla_{\nu}]f'(R) = \frac{8 \pi G}{c^4}T_{\mu \nu}
\end{equation}
For the simplest case $f(R)=R-2\Lambda$ we recover GR. For a FLRW metric, the Friedmann equations become (with $c=1$):
\begin{equation}
    3f'H^2 = 8\pi G \rho_m + \frac{1}{2}(f'R-f) - 3H \dot{f'}
\label{eqn:friedmann_fR}
\end{equation}
\begin{equation}
    -2f'\dot{H} = 8 \pi G (\rho_m + P_m) + \ddot{f'} - H \dot{f'}
\end{equation}
If one rewrites Equation \ref{eqn:friedmann_fR} to match what we get in GR $3H^2 = 8\pi G\rho$, we can identify an effective "curvature fluid" that under certain conditions 
would cause an accelerated Universe. This gives us an effective equation of state on $w_{eff}$:
\begin{equation}
    w_{eff} = -1 + \frac{1}{3f'H^2}\left[8\pi G(\rho_m + P_m) + \ddot{f'} - H \dot{f'} \right]
\end{equation}
Since acceleration is possible for $w<-\frac{1}{3}$ this means that a condition on any $f(R)$ theory would be:
\begin{equation}
    8\pi G(\rho_m + P_m) + \ddot{f'}-H\dot{f'}<2f'H^2
\end{equation}
$f(R)$ must also be able to reproduce GR during the early Universe and at high matter density regions to remain consistent with observational results. This screening effect is called 
the chameleon mechanism and the best-known $f(R)$ theories do possess it. \\

\noindent \textbf{\underline{Massive gravity}} \\
In quantum field theory, forces are mediated by particles, for instance photons for electromagnetism and gluons for the strong force. 
By analogy, gravity is expected to be mediated by a spin-2 particle called the graviton. The graviton is predicted to be a massless 
particle with an infinite interaction range, that can recover GR in a quantum field theory formalism. If we now assume that the graviton does in fact possess a mass, that would 
decrease the range of its gravitational interaction as well as its strength between distant objects (\cite{Fierz_1939}, \cite{deRham_2011}). The gravitational potential is then modified as follows:
\begin{equation}
    \Phi = - \frac{GM}{r}e^{-mr}
\end{equation}
where $r$ is the distance to an object, $M$ the mass of the object and $m$ the mass of the graviton. However, the introduction of massive gravitons produces multiple 
instabilities such as the inability to recover GR when $m$ goes to 0. It also has difficulty producing flat FLRW solutions.\\

\noindent \textbf{\underline{Breaking assumptions}} \\

The last class of models we will present here are the ones that break physical assumptions. Here we present an example of a model proposed by Gia Dvali, Gregory Gabadadze 
and Massimo Porrati (\cite{Dvali_2000}) where gravity can propagate in a 5-dimensional spacetime with an additional spatial dimension. We can think of it as our 4D world embedded 
in an infinite volume 5D space. The idea being that the gravity on large scales appears to be in 5D while on small scales in 4D, this is due to the Vainshtein effect. 
The action for such a theory looks like :
\begin{equation}
    S = 2M_5^3 \int d^5X \sqrt{-\tilde{g}}\tilde{R} + 2M_{pl}^2 \int d^4x \sqrt{-g}R + \int \sqrt{-g}\mathcal{L}_m + S_{GH-} + S_{GH+}
\end{equation}
where $M_5$ is the Planck mass $M_{pl}$ in 5D, $\tilde{g}$ and $\tilde{R}$ are the norm of the metric and the Ricci scalar in 5D. The $S_{GH-}$ and $S_{GH+}$ are Gibbons-Hawkings-York 
terms that serve as boundary terms commonly found in, but not limited to, quantum gravity theories.\\ 
The theory does predict a late time acceleration (\cite{Deffayet_2001}), however it suffers theoretical instabilities and unphysical solutions (\cite{Koyama_2007}). 
Observations have also ruled the theory out 
DGP theories (\cite{Fang_2008}). \\

